% =====================================================================
\chapter{Conclusion}
\label{sec:conclusion}
% =====================================================================
This project started from one question I could not let go of: if an examiner makes a mistake, who notices? The evidence
shows the question is real. Marking errors are documented, examiners are paid per script, most scripts are never seen by
a second person, and boards do not publish how often marks are wrong.

The prototype shows that a cheap, local audit layer can help. It reads students' handwriting well, its simple rules
catch arithmetic errors with certainty, and its AI second reader found a real six-mark under-marking in a dataset that
was described as verified. With the same re-checking effort, risk ranking found about 1.5x more planted errors than random
checking, and about 2x with the best combination of signals.

The prototype also shows clearly what does not work yet: reading the examiner's red ink. That is the first thing I would
fix, and it is exactly where an in-house GPU, larger models and fine-tuning on real marked scripts would make the biggest
difference. The AI should never give the final mark. Its job is to make sure that the scripts most likely to be wrong are
the ones a human looks at first.
